\documentclass[conference]{IEEEtran}
\usepackage{cite}
\usepackage{amsmath,amssymb,amsfonts}
\usepackage{graphicx}
\usepackage{textcomp}
\usepackage{xcolor}
\usepackage{booktabs}
\usepackage{array}

\begin{document}

\title{AgentDSS: A Human-in-the-Loop Sequential Decision Pipeline for Supply Chain Disruption Management}

\author{\IEEEauthorblockN{Meghashyam Sai Dontha}}

\maketitle

\begin{abstract}
Supply chain disruptions are typically managed through siloed, department-level decisions that ignore downstream consequences and give human managers no structured way to inject real-world knowledge that automated systems lack. This proposal presents AgentDSS, a human-in-the-loop sequential decision pipeline in which four role-specific large language model (LLM) agents (Supplier, Inventory, Logistics, Finance) generate decision options conditioned on all prior human decisions in the chain. At every step, a human may accept an agent's option or supply an unconstrained override, which is propagated as context to every subsequent agent. A deterministic Impact Preview heuristic, computed from a public shipment-level logistics dataset, estimates downstream consequences before a decision is confirmed. The proposal treats evaluation as a first-class deliverable: a stratified synthetic scenario suite, a controlled comparison between override-enabled and no-override baseline conditions scored by an independent judge model, pre-specified statistical testing with correction for multiple comparisons, a pilot validation phase, and explicit reproducibility and reporting standards.
\end{abstract}

\begin{IEEEkeywords}
decision support systems, large language model agents, human-in-the-loop systems, supply chain management, evaluation methodology, reproducibility
\end{IEEEkeywords}

\section{Introduction}
Modern supply chains are managed through department-level systems --- supplier, inventory, logistics, and finance --- that operate largely in isolation \cite{power2002}. When a disruption occurs, three structural problems recur. First, siloed decision-making: departments react independently, sometimes producing contradictory actions. Second, the absence of a human override mechanism: automated decision support systems (DSS) generate recommendations algorithmically with no structured channel for a manager's real-world knowledge to alter downstream computation. Third, blind decision-making: a manager acting at one point in the chain has no visibility into the consequences of that decision further downstream.

This proposal describes AgentDSS, a system in which each department is represented by a role-specific LLM agent. Agents generate options sequentially, each conditioned on the full history of prior human decisions, so that an override at any step changes the options generated by every subsequent agent. The system is deliberately framed as a human-in-the-loop sequential decision pipeline rather than a multi-agent system in the negotiation-based sense \cite{guo2024}: agents do not communicate with, or reach consensus with, one another. All coordination passes through a single shared decision log.

\section{Problem Statement}
\subsection{Siloed Decision-Making}
Departmental systems act on local information only, with no shared record of decisions made elsewhere in the chain during a disruption event.
\subsection{Absence of Human Override}
A manager who knows that stock can be sourced informally, or that a supplier relationship has off-the-record flexibility, has no way to inject that knowledge into an automated recommendation pipeline.
\subsection{Blind Downstream Consequences}
A decision made at one stage of the chain is committed without visibility into its effect on later stages, financially or operationally.

\section{Proposed Solution}
AgentDSS addresses each problem directly: shared context eliminates siloing, since every agent conditions on the complete decision history; free-text override at every step provides a channel for human knowledge; and a deterministic Impact Preview heuristic (Section~\ref{sec:preview}) estimates downstream consequences before a decision is confirmed.

\begin{table}[h]
\centering
\caption{Problem-to-Mechanism Mapping}
\begin{tabular}{p{3.8cm}p{3.8cm}}
\toprule
\textbf{Problem} & \textbf{Mechanism} \\
\midrule
Siloed decisions & Every agent receives the complete prior decision history \\
No human override & Free-text override accepted at every step; propagated immediately \\
Blind decision-making & Deterministic Impact Preview surfaces estimated downstream effects pre-confirmation \\
\bottomrule
\end{tabular}
\end{table}

\section{System Operation}
The pipeline executes as follows: (1) four agents are configured, each assigned a role and a partition of the backbone dataset; (2) the operator enters a disruption description; (3) Agent 1 reads its data partition and the disruption description and returns two to three ranked options, each with a confidence score; (4) the Impact Preview Engine computes an estimated downstream effect for each option using a precomputed heuristic rather than live agent calls; (5) the human accepts an option or supplies a free-text override; (6) the confirmed decision is written to the decision log and Context Manager; (7) Agent 2 reads the updated context and generates new options; (8) steps 3--7 repeat through Agent 4; (9) a complete decision log is produced.

\section{System Architecture}
The backend consists of four components: the \textbf{Agent Engine}, which reads dataset and context and calls the Claude API to return ranked options with confidence; the \textbf{Context Manager}, which maintains the decision chain and detects overrides by checking whether the chosen text matches a presented option; the \textbf{Impact Preview Engine}, which applies the heuristic below rather than live parallel LLM calls to remaining agents; and the \textbf{Decision Log Engine}, which persists every recommendation, confidence score, preview, and human decision to PostgreSQL.

\subsection{Impact Preview Heuristic}
\label{sec:preview}
The Impact score of a candidate decision $d$ is a weighted sum of three normalized deltas:
\begin{equation}
\text{Impact}(d) = w_c \Delta\text{cost}(d) + w_t \Delta\text{leadtime}(d) + w_r \Delta\text{risk}(d)
\end{equation}
where, for each quantity $x \in \{\text{cost}, \text{leadtime}, \text{risk}\}$,
\begin{equation}
\Delta x(d) = \frac{x(d) - x_0}{x_{\max} - x_{\min}}
\end{equation}
$x_0$ is the value in the scenario's baseline state, and $x_{\min}, x_{\max}$ are fixed bounds computed once from the full backbone dataset. Each delta lies in $[-1,1]$ and the score therefore lies in $[-1,1]$, where lower is better. The weights are fixed before evaluation at $w_c = 0.33$, $w_t = 0.33$, $w_r = 0.34$, and are reported alongside all results; equal weighting is adopted because no domain evidence favors one factor. The formulation is deterministic given the dataset and decision, whereas a live-LLM preview would introduce run-to-run variance that confounds the comparison in Section~\ref{sec:eval}.

\section{Backbone Dataset}
\label{sec:data}
The pipeline is grounded in a public Kaggle dataset of 5,000 shipment records with 14 fields: shipment ID, date, origin and destination port, transport mode, product category, distance (km), weight (MT), fuel price index, geopolitical risk score (0--10), weather condition, carrier reliability score (0.5--1.0), lead time (days), and a binary disruption flag \cite{kaggle}. Eight ports are represented, including Los Angeles. The dataset is used in full, without regional filtering.

\begin{table}[h]
\centering
\caption{Dataset Fields Used by Each Component}
\footnotesize
\begin{tabular}{p{2.4cm}p{2.1cm}p{2.9cm}}
\toprule
\textbf{Fields} & \textbf{Component} & \textbf{Role} \\
\midrule
Origin port, carrier reliability & Supplier agent & Candidate sourcing routes and carrier trust \\
Product category, weight & Inventory agent & Throughput proxy \\
Mode, distance, weather & Logistics agent & Route and mode alternatives \\
Fuel index, distance, weight & Finance agent; $\Delta$cost & Derived cost proxy \\
Lead time & All agents; $\Delta$leadtime & Time impact \\
Geopolitical risk & $\Delta$risk & Route risk \\
\bottomrule
\end{tabular}
\end{table}

The dataset has no native cost field, so cost is a derived proxy:
\begin{equation}
\text{cost} = P_{\text{fuel}} \cdot D_{\text{km}} \cdot W_{\text{MT}}
\end{equation}
with observed range $4.385\times10^{3}$ to $3.035\times10^{7}$. Fixed bounds for lead time are $[0.5, 236.39]$ days and for risk $[0, 10]$.

\subsection{Dataset Limitations}
The dataset contains no inventory levels, demand series, or financial statements. The Inventory agent therefore operates on weight and product category as a throughput proxy, and the Finance agent on the derived cost proxy; both are proxies, not measurements of inventory or financial position. The dataset is global rather than regional. Lead time is heavy-tailed (75th percentile about 21 days, maximum 236 days), so min-max scaling compresses most lead-time deltas toward zero. Several fields have near-uniform marginal distributions, which is consistent with generated data; the dataset's provenance and generation process are verified and reported in the final report, and results are presented as evidence about pipeline behavior on a realistic-format benchmark, not as claims about real-world supply chain performance.

\section{Decision Event Schema}
\label{sec:schema}
Each decision event is stored as a structured record so override behavior can be queried and scored programmatically: \texttt{event\_id}, \texttt{agent\_role}, \texttt{options\_presented} (with confidence scores), \texttt{decision\_type} (accepted\_option $\vert$ override), \texttt{decision\_text}, \texttt{override\_flag}, \texttt{preview\_shown}, \texttt{confidence\_score} (denormalized from the chosen option for cheap querying), and \texttt{timestamp}.

\section{Evaluation Methodology}
\label{sec:eval}
Evaluation is the central technical contribution of the project and is treated as a first-class deliverable.

\subsection{Pilot Validation}
Before full-scale evaluation, the complete pipeline is run end-to-end on a pilot set of 10 scenarios to confirm that judge scores, heuristic outputs, and override propagation behave as expected. Full-scale evaluation proceeds only after the pilot passes a manual sanity check; this ordering is a binding process constraint.

\subsection{Scenario Suite}
Sixty synthetic disruption scenarios are generated from the backbone dataset, stratified evenly across four templates at 15 scenarios per template: supplier delay, demand spike (modeled as a weight surge for a product category, since no demand series exists), logistics route failure, and combined multi-factor disruption. A scenario consists of a baseline state (an anchor shipment), a template, and disruption parameters (severity, affected node, time horizon) sampled across the range supported by the dataset. Each agent's candidate pool is the set of dataset records related to the affected node, and every pool contains at least four records. Parameter ranges and the sampling scheme are reported to allow independent verification of scenario diversity.

\subsection{Experimental Conditions}
Each scenario is run under two conditions: (1) a no-override baseline in which every agent's top-confidence option is accepted; and (2) an override condition driven by a scripted override policy. For scenario $s$, a pseudo-random generator seeded with the global seed and $s$ selects one step $k_s \in \{1,2,3,4\}$. At step $k_s$ the simulated operator rejects all presented options and supplies, as free text, the lowest-Impact alternative in the candidate pool that was not presented; at every other step the top-confidence option is accepted. Each scenario therefore contains exactly one deviation, satisfying the at-least-one requirement, and the deviation triggers the override branch of the Context Manager. Scripted overrides keep runs deterministic and repeatable and test the override mechanism, not the quality of human judgment. A small number of live human sessions are recorded for qualitative comparison and reported separately from the scripted-run statistics \cite{governance2026}.

\subsection{Scoring}
Each completed run is scored on: (1) a decision-quality rubric evaluated by an LLM-as-judge from a different provider family than the agent-generation model, with agreement against a human-rated subsample (size fixed in advance) reported as Cohen's kappa or Spearman correlation; (2) cost, lead-time, and risk deltas computed directly from the dataset fields; and (3) an override-acceptance pattern analysis, including cases where the top-confidence option carries a high Impact score.

\subsection{Ablation Studies}
\begin{itemize}
\item Does the Impact Preview measurably change the decision that follows it, compared to a run where the preview is suppressed?
\item Does agent confidence correlate with observed override rate?
\item Does context-conditioned option generation differ meaningfully from a static baseline in which each agent ignores prior decisions?
\end{itemize}

\subsection{Statistical Treatment}
The primary override-vs-baseline comparison and each ablation are evaluated with the Wilcoxon signed-rank test on paired scenario runs, chosen over a paired $t$-test because rubric scores are ordinal. Across the four planned comparisons (main comparison plus three ablations), a Bonferroni correction sets the significance threshold at $0.05/4 = 0.0125$. All test statistics, exact $p$-values, and effect sizes are reported in full.

\section{Reproducibility and Reporting Standards}
\label{sec:repro}
\subsection{Reproducibility}
\begin{itemize}
\item Scenario generation and override-step selection use a fixed random seed; the same seed reproduces the same 60 scenarios and the same overrides.
\item Agent API calls during evaluation use a fixed, low temperature.
\item Every prompt version is version-controlled and timestamped; any prompt revision after the pilot triggers a re-run of affected scenarios.
\item Impact Preview weights and normalization bounds are fixed in code and reported; they do not change between runs.
\item Raw output of every scenario run is stored, not only aggregate scores.
\end{itemize}

\subsection{Reporting}
\begin{itemize}
\item Results are reported as exact counts, test statistics, and $p$-values.
\item Ablation results that do not support a hypothesis are published alongside those that do.
\item API cost and end-to-end latency per scenario run are disclosed.
\end{itemize}

\section{Technology Stack}
Python 3 with pandas and NumPy for dataset processing; SciPy for statistical tests; SQLAlchemy with PostgreSQL for the decision log; pytest for unit tests; the Claude API for agent generation and a model from a different provider family for judge scoring; git and GitHub for version control. A FastAPI backend and React frontend are optional components for a live demonstration; the evaluation harness runs as a batch script and does not depend on them.

\section{Related Work and Novelty}
Existing work includes agentic consensus-seeking among LLM agents in inventory management \cite{jannelli2024}, standalone zero-shot inventory agents without inter-agent communication \cite{quan2024}, retailer-specific planning agents \cite{jd2025}, and multi-agent risk assessment for logistics hub selection \cite{zhao2026}. A separate literature on AI-assisted decision making studies when humans should rely on AI recommendations. Confidence scores can help calibrate trust, but calibrated trust alone does not guarantee better joint decisions \cite{zhang2020}; a confident AI prediction can still be wrong \cite{whotrust2023}; and reliance is commonly assessed by outcome, distinguishing over-reliance from under-reliance \cite{areyousure2024}. AgentDSS relates these findings to sequential supply chain decisions, where a decision's context propagates to later stages. None of the surveyed supply chain systems offers user-configurable agents, an unconstrained human override channel, a downstream consequence preview prior to commitment, or a published comparative evaluation of override-enabled operation against a no-override baseline with pre-specified statistical treatment and reproducibility guarantees.

\section{Preliminary Work and Project Plan}
\subsection{Preliminary Work}
Implemented and tested: the Decision Event schema and an in-memory Context Manager with automatic override detection, covered by unit tests; the Impact Preview heuristic evaluated against the backbone dataset; and a first Supplier agent that reads dataset records and ranks options by a deterministic composite confidence. This agent is a development stub and is replaced by Claude API agents.

\subsection{Phases}
The project is carried out by one person as sequential milestones.
\begin{itemize}
\item Phase 1: schema, Context Manager, Impact Preview, Supplier agent stub (in progress).
\item Phase 2: replace the stub with Claude API agents; connect Agent 2 and validate context-dependent option generation, the highest-risk assumption.
\item Phase 3: seeded scenario generator with four templates and candidate-pool filtering; finalize Impact Preview with baseline deltas.
\item Phase 4: Agents 3--4, database persistence of full decision logs, scripted override policy, baseline-vs-override runner, and the 10-scenario pilot.
\item Phase 5: judge integration and human-rated subsample with agreement statistics.
\item Phase 6: full 60-scenario evaluation, ablations, statistical analysis, reproducibility verification, and write-up; optional live demonstration.
\end{itemize}

\section{Limitations and Future Work}
Scripted overrides test the override mechanism, not the quality of human judgment; a user study is future work. The Inventory and Finance agents rely on proxies because the dataset lacks inventory and financial fields. LLM self-reported confidence may be miscalibrated relative to Impact Preview. A sensitivity analysis over Impact weights and robust scaling of heavy-tailed lead times is deferred. Adaptive agents that update confidence from accumulated override history are excluded from this version because they would make runs order-dependent and break the paired design; they are a direction for future work, as is a regional (for example, US West Coast) subset.

\section{Expected Outcomes}
\begin{itemize}
\item A working pipeline demonstrating the full agent cycle with human override at every step.
\item A comparative evaluation across the 60-scenario suite with pre-specified significance testing and effect sizes.
\item Ablation results for each core mechanism.
\item A complete, exportable decision log per scenario run.
\item A technical report meeting the reproducibility and reporting standards above.
\end{itemize}

\begin{thebibliography}{12}
\bibitem{jannelli2024} V. Jannelli et al., ``Agentic LLMs in the supply chain: Towards autonomous multi-agent consensus-seeking,'' arXiv:2411.10184, 2024.
\bibitem{quan2024} Y. Quan and Z. Liu, ``InvAgent: A large language model based multi-agent system for inventory management in supply chains,'' arXiv:2407.11384, 2024.
\bibitem{guo2024} T. Guo et al., ``Large language model based multi-agents: A survey of progress and challenges,'' arXiv:2402.01680, 2024.
\bibitem{jd2025} ``Leveraging LLM-based agents for intelligent supply chain planning,'' JD.com case study, arXiv:2509.03811, 2025.
\bibitem{zhao2026} M. Zhao et al., ``Leveraging large language models to enhance multi-agent risk assessment in supply chain networks (MARS framework),'' Int. J. Prod. Res., 2026.
\bibitem{power2002} D. J. Power, \emph{Decision Support Systems: Concepts and Resources for Managers}. Westport, CT: Quorum Books, 2002.
\bibitem{zhang2020} Y. Zhang, Q. V. Liao, and R. K. E. Bellamy, ``Effect of confidence and explanation on accuracy and trust calibration in AI-assisted decision making,'' in \emph{Proc. ACM FAccT}, 2020, doi:10.1145/3351095.3372852.
% TODO: add authors for the next three entries
\bibitem{whotrust2023} ``Who should I trust: AI or myself? Leveraging human and AI correctness likelihood to promote appropriate trust in AI-assisted decision-making,'' arXiv:2301.05809.
\bibitem{areyousure2024} ``Are you really sure? Understanding the effects of human self-confidence calibration in AI-assisted decision making,'' in \emph{Proc. ACM CHI}, 2024, doi:10.1145/3613904.3642671.
\bibitem{governance2026} ``Harnessing embodied agents: Runtime governance for policy-constrained execution,'' arXiv:2604.07833.
% TODO: fill uploader, title, URL, license, and access date from the Kaggle page
\bibitem{kaggle} ``Global supply chain risk and logistics (2024--2026),'' Kaggle dataset, 5,000 records.
\end{thebibliography}

\end{document}
